\documentclass[12pt]{article}
\usepackage{amsmath,amssymb,amsthm}
\usepackage[margin=1in]{geometry}

\title{Project Euler Problem 700: Eulercoin}
\author{}
\date{}

\begin{document}
\maketitle

\section{Problem Statement}
Consider the sequence $a_n = 1504170715041707 \cdot n \bmod 4503599627370517$.
An element $a_n$ is called an \emph{Eulercoin} if it is strictly smaller than all previously found Eulercoins (i.e., $a_n < a_k$ for all $k < n$ such that $a_k$ was an Eulercoin, and $a_n < a_j$ for all $j < n$).

Find the sum of all Eulercoins.

\section{Mathematical Analysis}

Let $N = 1504170715041707$ and $M = 4503599627370517$. Since $M$ is prime and $\gcd(N, M) = 1$, the sequence $\{a_n\}_{n=1}^{M-1}$ is a permutation of $\{1, 2, \ldots, M-1\}$.

An Eulercoin at index $n$ satisfies $a_n < a_j$ for all $1 \le j < n$. These are the \emph{record minima} (left-to-right minima) of the sequence.

\subsection{Two-Phase Algorithm}

\textbf{Phase 1 (Forward):} Iterate $a_n = (a_{n-1} + N) \bmod M$ and track minima. This finds Eulercoins with large values quickly.

\textbf{Phase 2 (Backward):} For small values $v = 1, 2, \ldots$, compute the index $n_v = v \cdot N^{-1} \bmod M$ where $N^{-1}$ is the modular inverse. Track decreasing indices to identify Eulercoins.

\subsection{Correctness}
A value $v$ is an Eulercoin if and only if its index $n_v$ is less than the index of any value $v' < v$ that is also an Eulercoin. By iterating upward from $v = 1$ and maintaining a running minimum of indices, we correctly identify all Eulercoins from the small-value side.

\section{Complexity Analysis}
\begin{itemize}
    \item Phase 1: $O(k)$ where $k$ is the index of the transition point ($\approx 10^7$).
    \item Phase 2: $O(V)$ where $V$ is the value of the last Phase 1 Eulercoin ($\approx 10^7$).
    \item Total: $O(10^7)$ time, $O(1)$ space.
\end{itemize}

\section{Answer}

\[
\boxed{1517926517777556}
\]

\end{document}
